% GENERATED FILE. Edit canonical lessons, then run npm run generate:books.
% canonical-source-hash: fnv1a64:2259e3eca7460d94
% canonical-lessons: ZH-C362-dasheng, ZH-C362-bixu, ZH-W362-xu-must, ZH-C362-luan, ZH-C362-gongzheng, ZH-C362-yiwu

\chapter{What You Have to Do}
\label{ch:zh-what-you-have-to-do}
\hlchaptermodality{362}

\begin{chapteropening}
\textbf{By the end of this chapter:} \emph{I can say what you have to do, ask someone not to be loud, and say what is fair.}

Complete the last lesson of chapter 362: I can say what you have to do, ask someone not to be loud, and say what is fair.
\end{chapteropening}

\section[dàshēng]{dàshēng — loudly; in a loud voice}

\label{lesson:ZH-C362-dasheng}



\begin{quote}
Before the new one: say the Mandarin for to propose a toast, then the Mandarin for to look down on someone.
\end{quote}

\subsection*{You'll want to know: dàshēng}
\emph{dàshēng} — \textquotedblleft{}loudly; in a loud voice\textquotedblright{}.

\begin{quote}
\emph{qǐng búyào dàshēng shuōhuà} — Please don't talk loudly.
\end{quote}

Its characters come later; for now it is a word for the ear.

\begin{cousinweb}
\emph{dàshēng} is built from \zh{大} 'big' and \emph{shēng}, both inherited from Middle Chinese, \emph{dajH} and \emph{syeng} (Guangyun, in Baxter's transcription), which glosses the first 'big' and the second 'a sound'.
\end{cousinweb}

\begin{grammarlens}[title={What you've built}]
One more word for giving reasons.
\end{grammarlens}

\subsection*{Your turn}
\begin{itemize}
\raggedright
  \item \emph{Say it:} \emph{dàshēng}
  \item \emph{Say it:} \emph{dàshēng}, once more
  \item \emph{Say it:} \emph{qǐng búyào dàshēng shuōhuà} — Please don't talk loudly.
  \item \emph{Recall:} say the Mandarin for to propose a toast, then the Mandarin for to look down on someone, then say \emph{dàshēng} again
\end{itemize}

\begin{tcolorbox}[breakable,colback=teal!4,colframe=teal!35!black,title={Before you move on}]
What does dàshēng mean? (\textquotedblleft{}Loudly; in a loud voice\textquotedblright{}.) Say it once more.
\end{tcolorbox}
\section[\texorpdfstring{\zh{必须} (bìxū)}{必须 (bìxū)}]{\zh{必须} (bìxū) — must; it is required that}

\label{lesson:ZH-C362-bixu}



\begin{quote}
Before the new one: say the Mandarin for to look down on someone, then the Mandarin for in a loud voice.
\end{quote}

\subsection*{You'll want to know: \zh{必须}}
\textbf{\zh{必须}} — \emph{bìxū} — \textquotedblleft{}must; it is required that\textquotedblright{}.

\begin{quote}
\textbf{\zh{上课的时候必须关手机}} \emph{shàngkè de shíhou bìxū guān shǒujī} — You have to switch your phone off in class.
\end{quote}

The new character, in pieces: \zh{须} — nine strokes in two pieces: on the left three left-falling strokes, one under another; on the right the page piece of \zh{题}, six strokes. The next lesson puts it in your hand.

\begin{cousinweb}
\textbf{\zh{必须}} (\emph{bìxū}) is built from \zh{必} 'surely' and \emph{xū}, both inherited from Middle Chinese, \emph{pjit} and \emph{sju} (Guangyun, in Baxter's transcription), which glosses the first 'certain' and the second 'what the mind wants' and, quoting an older dictionary, 'the hair of the face'.
\end{cousinweb}

\begin{grammarlens}[title={What you've built}]
One more word for giving reasons.
\end{grammarlens}

\subsection*{Your turn}
\begin{itemize}
\raggedright
  \item \emph{Say it:} \emph{bìxū}
  \item \emph{Say it:} \emph{bìxū}, once more
  \item \emph{Say it:} \emph{shàngkè de shíhou bìxū guān shǒujī} — You have to switch your phone off in class.
  \item \emph{Recall:} say the Mandarin for to look down on someone, then the Mandarin for in a loud voice, then say \emph{bìxū} again
\end{itemize}

\begin{tcolorbox}[breakable,colback=teal!4,colframe=teal!35!black,title={Before you move on}]
What does \zh{必须} mean? (\textquotedblleft{}Must; it is required that\textquotedblright{}.) Say it once more.
\end{tcolorbox}
\section[\texorpdfstring{\zh{须} (xū)}{须 (xū)}]{\zh{须} — three sweeps, then the page piece}

\label{lesson:ZH-W362-xu-must}



\begin{quote}
The word first, then the character inside it.

\begin{itemize}
\raggedright
  \item \emph{Recall:} say \emph{bìxū} — \textbf{R1}, one lesson back, before it decays
\end{itemize}

\begin{quote}
\zh{题} \zh{必}
\end{quote}

\emph{Read it:} each one aloud, with nothing said first — \emph{tí}, \emph{bì}
\end{quote}

\subsection*{Script — three sweeps, then the page piece}
\hlblockfigure{\detokenize{figures/ZH-W362-xu-must-filmstrip.pdf}}{How \zh{须} is written, stroke by stroke}

\begin{quote}
\zh{须}
\end{quote}

\textbf{Two pieces, left before right.} Count the pieces, not the strokes.

\textbf{The left piece, three strokes}: three short strokes falling to the left, one under another, the last the longest.

\textbf{The right piece, the page piece, six strokes}, as on the right of \zh{题}: a horizontal, a short left-falling stroke, a vertical, a horizontal that turns down, a left-falling stroke inside, then a dot at the foot.

3 + 6 = nine.

\subsection*{Your turn}
\begin{itemize}
\raggedright
  \item \emph{Write it:} \zh{须}, all nine, with the model in front of you
  \item \emph{Cover:} write it again from memory, and count the pieces, not the strokes
\end{itemize}

\begin{tcolorbox}[breakable,colback=teal!4,colframe=teal!35!black,title={Before you move on}]
How many pieces in \zh{须}? (\textbf{Two}.) Where have you written the right piece? (\textbf{On the right of \zh{题}}.) How many strokes on the left? (\textbf{Three, all falling to the left}.)

Source: \href{https://raw.githubusercontent.com/chanind/hanzi-writer-data/68d10a4b21150cae5e1ebbd223eed289cf32d90c/data/\%E9\%A1\%BB.json}{Hanzi Writer Data \zh{须}, pinned stroke paths and medians}.
\end{tcolorbox}
\section[luàn]{luàn — messy; in a mess}

\label{lesson:ZH-C362-luan}



\begin{quote}
Before the new one: say the Mandarin for in a loud voice, then the Mandarin for must (a firm requirement).
\end{quote}

\subsection*{You'll want to know: luàn}
\emph{luàn} — \textquotedblleft{}messy; in a mess\textquotedblright{}.

\begin{quote}
\emph{nǐ de fángjiān tài luàn le} — Your room is far too messy.
\end{quote}

Its characters come later; for now it is a word for the ear.

\begin{cousinweb}
\emph{luàn} is inherited from Middle Chinese \emph{lwanH} (Guangyun, in Baxter's transcription), which glosses it both 'to put in order' and 'disorder'.
\end{cousinweb}

\begin{grammarlens}[title={What you've built}]
One more word for giving reasons.
\end{grammarlens}

\subsection*{Your turn}
\begin{itemize}
\raggedright
  \item \emph{Say it:} \emph{luàn}
  \item \emph{Say it:} \emph{luàn}, once more
  \item \emph{Say it:} \emph{nǐ de fángjiān tài luàn le} — Your room is far too messy.
  \item \emph{Recall:} say the Mandarin for in a loud voice, then the Mandarin for must (a firm requirement), then say \emph{luàn} again
\end{itemize}

\begin{tcolorbox}[breakable,colback=teal!4,colframe=teal!35!black,title={Before you move on}]
What does luàn mean? (\textquotedblleft{}Messy; in a mess\textquotedblright{}.) Say it once more.
\end{tcolorbox}
\section[\texorpdfstring{\zh{公正} (gōngzhèng)}{公正 (gōngzhèng)}]{\zh{公正} (gōngzhèng) — fair; impartial}

\label{lesson:ZH-C362-gongzheng}



\begin{quote}
Before the new one: say the Mandarin for must (a firm requirement), then the Mandarin for in a mess.
\end{quote}

\subsection*{You'll want to know: \zh{公正}}
\textbf{\zh{公正}} — \emph{gōngzhèng} — \textquotedblleft{}fair; impartial\textquotedblright{}.

\begin{quote}
\textbf{\zh{我觉得校长很公正}} \emph{wǒ juéde xiàozhǎng hěn gōngzhèng} — I think the head is very fair.
\end{quote}

\begin{cousinweb}
\textbf{\zh{公正}} (\emph{gōngzhèng}) is built from \zh{公} and \zh{正} 'proper', both inherited from Middle Chinese, \emph{kuwng} and \emph{tsyengH} (Guangyun, in Baxter's transcription), which glosses the first 'upright' and 'selfless', and the second 'proper'.
\end{cousinweb}

\begin{grammarlens}[title={What you've built}]
One more word for giving reasons.
\end{grammarlens}

\subsection*{Your turn}
\begin{itemize}
\raggedright
  \item \emph{Say it:} \emph{gōngzhèng}
  \item \emph{Say it:} \emph{gōngzhèng}, once more
  \item \emph{Say it:} \emph{wǒ juéde xiàozhǎng hěn gōngzhèng} — I think the head is very fair.
  \item \emph{Recall:} say the Mandarin for must (a firm requirement), then the Mandarin for in a mess, then say \emph{gōngzhèng} again
\end{itemize}

\begin{tcolorbox}[breakable,colback=teal!4,colframe=teal!35!black,title={Before you move on}]
What does \zh{公正} mean? (\textquotedblleft{}Fair; impartial\textquotedblright{}.) Say it once more.
\end{tcolorbox}
\section[yìwù]{yìwù — duty; obligation}

\label{lesson:ZH-C362-yiwu}



\begin{quote}
Before the new one: say the Mandarin for in a mess, then the Mandarin for impartial.
\end{quote}

\subsection*{You'll want to know: yìwù}
\emph{yìwù} — \textquotedblleft{}duty; obligation\textquotedblright{}.

\begin{quote}
\emph{wǒ juéde zhè shì dàjiā de yìwù} — I think it's everyone's duty.
\end{quote}

Its characters come later; for now it is a word for the ear.

\begin{cousinweb}
\emph{yìwù} is built from \emph{yì} and \emph{wù}, both inherited from Middle Chinese, \emph{ngjeH} and \emph{mjuH} (Guangyun, in Baxter's transcription), which glosses the first 'benevolence and right conduct' and the second 'affairs'.
\end{cousinweb}

\begin{grammarlens}[title={What you've built}]
One more word for giving reasons.
\end{grammarlens}

\subsection*{Your turn}
\begin{itemize}
\raggedright
  \item \emph{Say it:} \emph{yìwù}
  \item \emph{Say it:} \emph{yìwù}, once more
  \item \emph{Say it:} \emph{wǒ juéde zhè shì dàjiā de yìwù} — I think it's everyone's duty.
  \item \emph{Recall:} say the Mandarin for in a mess, then the Mandarin for impartial, then say \emph{yìwù} again
\end{itemize}

\begin{tcolorbox}[breakable,colback=teal!4,colframe=teal!35!black,title={Before you move on}]
What does yìwù mean? (\textquotedblleft{}Duty; obligation\textquotedblright{}.) Say it once more.
\end{tcolorbox}
